\documentclass{article}
\usepackage[margin=1in]{geometry} % margins
\usepackage[justification=centering]{caption}
%\usepackage{calrsfs}
\usepackage{amsmath}
\usepackage{amsfonts}
\usepackage{amssymb}
\usepackage{parskip}
\usepackage{graphicx}
\usepackage{enumerate}
\usepackage{enumitem}
\usepackage{amsthm}
\usepackage{float}
\usepackage{tikz-cd}
\usepackage{ytableau}
\usepackage{dsfont}

\renewcommand{\implies}{\Rightarrow}
\newcommand{\longimplies}{\Longrightarrow}
\newcommand{\N}{\mathbb{N}}
\newcommand{\Z}{\mathbb{Z}}
\newcommand{\Q}{\mathbb{Q}}
\newcommand{\R}{\mathbb{R}}
\newcommand{\C}{\mathbb{C}}
\newcommand{\D}{\mathbb{D}}
\renewcommand{\H}{\mathbb{H}}
\newcommand{\F}{\mathbb{F}}
\newcommand{\K}{\mathbb{K}}
\newcommand{\Lf}{\mathbb{L}}
\newcommand{\Sp}{\mathbb{S}}
\newcommand{\cl}[1]{\overline{#1}}
\renewcommand{\emptyset}{\varnothing}
\newcommand{\norm}[1]{\left\|#1\right\|}
\newcommand{\maxnorm}[1]{\left\|#1\right\|_{\text{max}}}
\newcommand{\opnorm}[1]{\left\|#1\right\|_{\mathrm{op}}}
\newcommand{\supnorm}[1]{\left\|#1\right\|_{\infty}}
\newcommand{\id}{\operatorname{id}}
\newcommand{\sspan}{\operatorname{span}}
\newcommand{\tr}{\operatorname{tr}}
\newcommand{\ch}{\operatorname{char}}
\newcommand{\poly}{\mathrm{P}}
\newcommand{\mtx}{\mathcal{M}}
\newcommand{\seq}{\mathbf{Seq}}
\newcommand{\vv}[1]{\mathbf{#1}}
\newcommand{\dirsum}{\oplus}
\newcommand{\rank}{\operatorname{rank}}
\newcommand{\im}{\operatorname{im}}

\newcommand{\cj}[1]{\overline{#1}}
\newcommand{\inprod}[2]{\left\langle#1,#2\right\rangle}
\renewcommand{\Re}{\operatorname{Re}}
\renewcommand{\Im}{\operatorname{Im}}

\newcommand{\res}{\operatorname{res}}

\newcommand{\floor}[1]{\left\lfloor#1\right\rfloor}

% Theorem environments
\usepackage{amsthm}

\theoremstyle{definition} % The "plain" style italicizes all body text.
	\newtheorem{thm}{Theorem}
		\numberwithin{thm}{section} % Theorem numbers are determined by section.
	\newtheorem{lemma}[thm]{Lemma}
	\newtheorem{prop}[thm]{Proposition}
	\newtheorem{cor}[thm]{Corollary}
	\newtheorem{conj}[thm]{Conjecture}

\theoremstyle{definition} % The "definition" style does not italicize body text..
    \newtheorem{defn}[thm]{Definition}
	\newtheorem{example}[thm]{Example}
    \newtheorem{exercise}[thm]{Exercise}

\usepackage{mathtools}

\usetikzlibrary{decorations.markings,positioning,decorations.text}

\def\xr{5}
\def\yr{5}

%\usepackage{titlesec}
%\titleformat{\section}[hang]{\normalfont\fontsize{14}{0}\bfseries}{Problem 8.\thesection\ }{0pt}{}

\begin{document}

\section{Week 10}

\begin{defn}
    If $E \in \mathcal{E}, f: E \to \cl{\R}$ then we say that $f$ is Lebesgue measurable if $f^{-1}(B)$ is Lebesgue measurable for all $B \in \mathcal{B}(\cl{\R})$.
\end{defn}

\begin{lemma}
    If $f: E \to \cl{\R}$ and if we define
    \[
        \tilde{f}(x) = \begin{cases}
            f(x) & x \in E \\
            0 & x \notin E
        \end{cases}  
    \]
    then $f$ is Lebesgue measurable if and only if $\tilde{f}$ is Lebesgue measurable.
\end{lemma}
\begin{proof}
If $B \in \mathcal{B}(\cl{\R})$ then if $0 \notin B$, $f^{-1}(B) = \tilde{f}^{-1}(B)$ and if $0 \in B$ then $\tilde{f}^{-1}(B) = f^{-1}(B) \cup E^c$, hence $f^{-1}(B) = f^{-1}(B) \cap E$.
\end{proof}

\begin{defn}
    We define $\mathds{1}_E(x) = \begin{cases}
        1 & x \in E \\ 0 & x \notin E
    \end{cases}$. If it is Lebesgue measurable then so is $E$.

    If $E \in \mathcal{E}$ we say that $\varphi: E \to \cl{\R}$ is simple if $\varphi = \sum_{k = 1}^n \alpha_k \mathds{1}_{E_k}$ where $E_k \subseteq E$ are Lebesgue measurable and $\alpha_k \in \R$. 
    
    Note that you can always rewrite the representative so that the $E_k$ are disjoint and the $\alpha_k$'s are unique. This representatation is called canonical.
\end{defn}

\begin{lemma}
    Let $f: E \to \R$ be a bounded Lebesgue measurable function with $E \in \mathcal{E}$, and let $\varepsilon > 0$. Then there exist simple functions $\varphi_\varepsilon$ and $\psi_\varepsilon$ such that $\varphi_\varepsilon \leq f \leq \psi_\varepsilon$ and $|\psi_\varepsilon(x) - \varphi_{\varepsilon}(x)| < \varepsilon$ for all $x \in E$. 
\end{lemma}
\begin{proof}
Suppose that $M > |f(x)|$ for all $x \in E$. Let $y_i \in \R$ be an increasing sequence satisfying $y_1 = -M$ and $y_n = M$. Then we let $E_k = f^{-1}((y_k, y_{k + 1}))$ for all $1 \leq k \leq n - 1$. We set $\varphi = \sum_{k = 1}^{n - 1} y_k\mathds{1}_{E_k}$ and $\psi = \sum_{k = 1}^{n - 1} y_{k + 1}\mathds{1}_{E_k}$. By definition $\varphi \leq f \leq \psi$ and $|\varphi - \psi| \leq \max_k|y_{k + 1} - y_k|$. By choosing $y_k$ to be close enough, we get the second desired condition.
\end{proof}

\begin{thm}
Let $f: E \to \cl{\R}$ be Lebesgue measurable. Then 
\begin{enumerate}[label = (\roman*)]
    \item If $f \geq 0$ then there exists an increasing sequence of simple functions $\varphi_n \geq 0$ such that $\lim_{n \to \infty} \varphi(x) = f(x)$ for all $x \in E$.
    \item In general, there exists a sequence of simple functions $\varphi_n$ such that $\lim_{n \to \infty} \varphi_n(x) = f(x)$ and $|\varphi_n(x)| \leq |f(x)|$ for all $x \in E$.
\end{enumerate}
\end{thm}
\begin{proof}
(ii) follows from (i) by writing $f = f_+ - f_-$ and finding simple $\psi_n, \chi_n$ that approximate $f_+, f_-$ in the way of (i) and defining $\varphi_n = \psi_n - \chi_n$. Since at most one of $f_+, f_-$ is non-zero at a given $x$, the same is the case for $\psi_n, \chi_n$ and hence the absolute value condition holds. The pointwise convergence obviously holds as well.

For (i), if $f \geq 0$, if we set $f_n = \max\{f, n\}$ then by the previous lemma, there is an increasing sequence $\psi_n \leq f_n$ and $|f_n - \psi_n| < 1/n$ for all $n$. Then we take $\varphi_1 = \max\{\psi_1, 0\}$ and $\varphi_k = \max\{0, \psi_k, \varphi_{k - 1}\}$ for $k \geq 2$. By induction, we find that $\varphi_k$'s are step functions such that $0 \leq \varphi_k \leq \varphi_{k + 1} \leq f_k$, and $|f_k - \varphi_k| < 1/k$. With this we get $\varphi_n(x) \to f(x)$ for all $x \in E$.
\end{proof}

\begin{thm}[Egorov]
    Let $f_n: E \to \R$ be a sequence of Lebesgue measurable functions converging pointwise to $f: E \to \R$. We also assume that $m(E) < \infty$. Then for all $\varepsilon > 0$, there is a closed set $F \subseteq E$ such that $m(E \setminus F) < \varepsilon$ and $f_n \to f$ converges uniformly on $F$ (that is, for all $\delta > 0$ there is some $N$ such that $|f(x) - f_n(x)| < \delta$ for all $x \in F$ and $n \geq N$).
\end{thm}

\begin{lemma}
    In the setting of the theorem, for all $u > 0$ and $\delta > 0$ there is some $A \subseteq E$ and an $N \geq 0$ such that $m(E \setminus A) < \delta$ and for all $x \in A$ and $n \geq N$ we have $|f_n(x) - f(x)| < u$.
\end{lemma}
\begin{proof}
Set $E_n = \{x \in E: |f_n(x) - f(x)| < u, \forall k \geq n\} = \cap_{k = n}^\infty \{x: |f_k(x) - f(x)| < u\}$. Since $f_n \to f$ pointwise, $E = \cup_{n = 1}^\infty E_n$. Furthermore, $E_n$ are increasing ($E_n \subseteq E_{n + 1}$ for all $n$). We then have $m(E) = \lim_{n \to \infty} m(E_n)$. Since $m(E) < \infty$, there is some $N$ such that $|m(E) - m(E_N)| < \delta$, and hence $m(E \setminus E_N) = m(E) - m(E_N) < \delta$. Taking $A = E_N$ completes the proof.
\end{proof}

\begin{proof}[Proof (Egorov)]
By the lemma, for all $n$ there is some $A_n$ such that $m(E \setminus A_n) < \varepsilon / 2^n$ and there exists $M(n)$ such that $|f(x) - f_k(x)| < 1/n$ for all $k \geq M(n)$. Setting $A = \cup_{n = 1}^\infty A_n$, we have
\begin{align*}
    m(E \setminus A) 
    &= m(E \cap A^c) \\ 
    &= m\left(E \cap \left(\bigcup_{n = 1}^\infty A_n^c\right)\right) \\
    &= m\left(\bigcup_{n = 1}^\infty E \setminus A_n\right) \\
    &\leq \sum_{n = 1}^\infty m(E \setminus A_n) \\
    &\leq \sum_{n = 1}^\infty \frac{\varepsilon}{2^n} \\
    &= \varepsilon
\end{align*}
Now if $x \in A$ then for all $n$, $x \in A_n$, so for all $k \geq M(n)$, we have $|f_k(x) - f(x)| < 1/n$. We thus see that $f_k \to f$ uniformly on $A$.

We now take $F$ closed such that $m(A \setminus F) \leq \varepsilon$ with $F \subseteq A$ to get $m(E \setminus F) \leq m(E \setminus A) + m(A \setminus F) \leq 2\varepsilon$.
\end{proof}

\begin{thm}[Lusin]
    Let $f: E \to \R$ be Lebesgue measurable with $m(E) < \infty$. Then for all $\varepsilon > 0$ there is a closed set $F$ and a continuous $g: F \to \R$ such that $m(E \setminus F) < \varepsilon$ and $f = g$ on $F$.
\end{thm}

\begin{prop}
    Let $\varphi: E \to \R$ be a simple function and let $\varepsilon > 0$. Then there exists a closed $F \subseteq E$ and a continuous $g: \R \to \R$ such that $m(E \setminus F) < \varepsilon$ and $g = \varphi$ on $F$.
\end{prop}
\begin{proof}
Write $\varphi = \sum_{k = 1}^\infty \alpha_k \mathds{1}_{E_k}$ with $E_k$ being disjoint. Take $F_k$ closed and bounded such that $F_k \subseteq E_k$ with $m(E_k \setminus F_k) < \varepsilon / n$. Since $F_k$'s are disjoint compact sets, we have $\delta = \min_{i \neq j}\inf_{x \in F_i, y \in F_j} |x - y| > 0$. For any set $A$, define $d(x, A) = \inf_{y \in A}|x - y|$.

\begin{exercise}
    Show that $x \mapsto d(x, A)$ is continuous.
\end{exercise}
\begin{proof}
    TODO
\end{proof}

Let $\rho(x): \R_{\geq 0} \to \R$ as $1$ on $[0, \delta/3]$, 0 on $[2\delta/3, \delta]$, and linear on $[\delta/3, 2\delta/3]$ such that $\rho$ becomes connected. Define $g_i = \alpha_i \rho(d(x, F))$. Each $g_i$ is continuous and if $x \in F_i$ then $\rho(d(x, F_i)) = 1$ and if $x \in F_j$ for $j \neq i$ then $\rho(d(x, F_j)) = 0$ because $d(x, F_i) \geq \delta$. We then define $g(x) = \sum_{i = 1}^n g_i(x)$.

Define $F = \bigcup_{n = 1}^\infty F_n$, and note that $g(x) = \varphi(x)$ for $x \in F$ and $m(E \setminus F) \leq \sum_{k = 1}^n m(E_k \setminus F_k) \leq n \varepsilon / n = \varepsilon$.

\begin{lemma}
    If $E$ has finite measure then for any $\varepsilon > 0$ there is some compact $K$ such that $m(E \setminus K) < \varepsilon$.
\end{lemma}
\begin{proof}
Set $E_n = E \cap [-n, n]$. Then since $\bigcup_n E_n = E$ and $m(E) = \lim_{n \to \infty} m(E_n)$, we can find $E_k$ such that $m(E \setminus E_k) < \varepsilon$.
\end{proof}

Then take $F \subseteq E_k$ with $F$ closed and $m(E_k \setminus F) < \varepsilon \implies m(E \setminus F) \leq m(E \setminus E_k) - m(E_k \setminus F_k) \leq 2 \varepsilon$.
\end{proof}

\begin{proof}[Proof (Lusin)]
Take $\varphi_n$ simple functions such that $\varphi$ converge to $f$ pointwise and let $g_n$ be a sequence of continuous functions such that $g_n = \varphi_n$ on a closed set $F_k$ with $m(E \setminus F_n) \leq \varepsilon / 2^n$. By Egorov, there is a closed $G \subseteq E$ with $m(E \setminus G) < \varepsilon$ such that $\varphi_n \to f$ uniformly on $G$. Take $F = \left(\bigcap_{n = 1}^\infty F_n\right) \cap G$. Then $F$ is closed and $m(E \setminus F) \subseteq m(E \setminus G) + \sum_n m(E \setminus F_n) \leq 2\varepsilon$. On $F$, we have $g_n = \varphi_n$ for all $n$, so $g_n \to f$ uniformly on $F$. Thus the funciton $g(x) = \lim_{n \to \infty} g_n(x)$ for $x \in F$ in continuous.
\end{proof}

\section{Week 11}

\begin{defn}
We define \[ \int_E \sum_{j = 1}^n \alpha_j \mathds{1}_{E_j} = \sum_{j = 1}^n \alpha_j m(E_j) \]
and 
\[ \int_E f = \sup\left\{\int_E \varphi: \varphi \leq f, \varphi \text{ simple} \right\} \]
\end{defn}

\begin{thm}[Monotone Convergence Theorem]
Let $f_n: E \to [0, \infty]$ be measurable, $f_n \leq f_{n + 1}$ for all $n$. Let $f = \lim_{n \to \infty} f_n$ pointwise. Then $\lim_{n\to \infty} \int_E f_n = \int_E f$.
\end{thm}
\begin{proof}
Since $f_n \leq f$, we have $\int_E f_n \leq \int_E f \implies \lim_{n\to\infty}\int_E f_n \leq \int_E f$. Let $0 < \alpha < 1$ and let $\varphi: E \to [0, \infty)$ be simple with $\varphi \leq f$. Suffices to show that $\lim_{n \to \infty} \int_E f_n \geq \alpha\int_E \varphi$. Let $E_k = \{x \in E: f_k(x) \geq \alpha \varphi(x)\}$ and note that $E_k \subseteq E_{k + 1}$ and $E = \bigcup_{k = 1}^\infty E_k$. We then have
\[ \lim_{n \to \infty} \int_E f_n \geq \int_E f_k \geq \int_E \mathds{1}_{E_k} f_k \geq \int_E \mathds{1}_{E_k} \alpha \varphi = \alpha\int_E \mathds{1}_{E_k} \varphi \]
Writing $\varphi = \sum_{j = 1}^n \alpha_j\mathds{1}_{F_j}$, we have
\[ \alpha\int_E \mathds{1}_{E_k} \varphi = \alpha\sum_{j = 1}^n \alpha_j m(F_j \cap E_k) \leq \lim_{n \to \infty} \int_E f_n \]
This holds for all $k$ so take $k \to \infty$. Since $F_j \cap E_k \subseteq F_j \cap E_{k + 1}$ for all $k$ and $\bigcup_{k = 1}^\infty F_j \cap E_k = F_j$ so by continuity of measure $\lim_{n \to \infty} m(F_j \cap E_n) = m(F_j)$. Thus we have $\lim_{n \to \infty} \int_E f_n \geq \alpha \sum_{j = 1}^n \alpha_j m(F_j) = \alpha \int_E \varphi$. Taking $\alpha \to 1$ and using the definiion of Lebesgue integral, we see that $\lim_{n \to \infty} \int_E f_n \geq \int_E f$. Hence the equality holds.
\end{proof}

\begin{lemma}
If $f, g: E \to [0, \infty]$ are measurable, then $\int_E f + g = \int_E f + \int_E g$.
\end{lemma}
\begin{proof}
Since $f, g$ are nonnegative and measurable, there exist nonnegative simple functions $\varphi_n, \psi_n$ that are increasing to $f, g$ and $\varphi_n \to f, \psi_n \to g$ pointwise. Then 
\[ \int_E f + g = \int_E \lim_{n \to \infty} (\phi_n + \psi_n) = \lim_{n \to \infty} \int_E \phi_n + \psi_n = \lim_{n \to \infty} \int_E \phi_n + \lim_{n \to \infty} \int_E \psi_n = \int_E f + \int_E g \]
\end{proof}

\begin{cor}
If $f_n: E \to [0, \infty]$ are measurable then $\int_E \sum_{n = 1}^\infty f_n = \sum_{n = 1}^\infty \int_E f_n$
\end{cor}
\begin{proof}
Just apply MCT.
\end{proof}

\begin{prop}
Let $f: E \to [0, \infty]$ be measurable. Then $\int_E f = 0$ iff $f = 0$ a.e. on $E$.
\end{prop}
\begin{proof}
If $f = 0$ a.e. on $E$ then if $\varphi = \sum_{j = 1}^n \alpha_j \mathds{1}_{E_j}$ with $E_j$ disjoint, $\alpha_j \geq 0$, satisfies $\varphi \leq f$, then if $\alpha_j > 0$, we have $m(E_j) = 0$. Thus $0 \leq \int_E f \leq \int_E \varphi \leq 0$ so $\int_E f = 0$.

If $\int_E f = 0$, consider the set $f^{-1}((0, \infty]) = \bigcup_{k = 1}^\infty f^{-1}((1/k, \infty])$, and define $E_k := f^{-1}((1/k, \infty])$. We then have 
\[ 0 = \int_E f \geq \int_E \mathds{1}_{E_k} f \geq \int_E \frac{1}{k}\mathds{1}_{E_k} = \frac{m(E_k)}{k} \geq 0 \]
Thus $m(E_k) = 0$ so $m(f^{-1}((0, \infty])) = \lim_{n \to \infty} m(E_k) = 0$. Since the set on which $f > 0$ is of measure zero, we have $f = 0$ a.e. on $E$.
\end{proof}

\begin{thm}[Better MCT]
Let $f_n, f: E \to [0, \infty]$ be measurable. Assume that $f_n \leq f_{n + 1}$ a.e. and assume $f = \lim_{n \to \infty} f_n$ a.e. on $E$. Then $\lim_{n \to \infty} \int_E f_n = \int_E f$
\end{thm}
\begin{proof}
Let $E_0 = \{x \in E: \lim_{n \to \infty} f_n(x) = f(x), f_n(x) \leq f_{n + 1}(x)\}$. Then $m(E \setminus E_0) = 0$ ($E \setminus E_0$ is a union of sets of measure zero). Now, we have 
\[ \lim_{n \to \infty} \int_E f_n = \lim_{n \to \infty} \int_E f_n (\mathds{1}_{E_0} + \mathds{1}_{E \setminus E_0}) = \lim_{n \to \infty} \int_E f_n \mathds{1}_{E_0} + \int_E f_n \mathds{1}_{E \setminus E_0} = \int_E \lim_{n \to \infty} f_n \mathds{1}_{E_0} = \int_{E_0} f = \int_E f \]
\end{proof}

\begin{thm}[Fatou's lemma]
If $f_n: E \to [0, \infty]$ are measurable then 
\[ \int_E \liminf_{n \to \infty} f_n \leq \liminf_{n \to \infty} \int_E f_n. \]
In particular, if $f_n \to f$ a.e. then $\int_E f \leq \liminf_{n \to \infty} \int_E f_n$.
\end{thm}
\begin{proof}
\[ \int_E \liminf_{n \to \infty} f_n = \int_E \lim_{n \to \infty}\left(\inf_{k \geq n} f_k\right) = \lim_{n \to \infty} \int_E \inf_{k \geq n} f_k \leq \lim_{n \to \infty}\inf_{k \geq n} \int_E f_k = \liminf_{n \to \infty} \int_E f_n \]
\end{proof}

\begin{thm}[Markov's inequality]
If $f: E \to [0, \infty]$ is measurable then for all $\alpha > 0$, $m(f^{-1}([\alpha, \infty])) \leq \frac{1}{\alpha} \int_E f$.
\end{thm}
\begin{proof}
Let $F = f^{-1}([\alpha, \infty])$, then $f(x)/\alpha \geq 1$ on $F$, so 
\[ m(F) = \int_E \mathds{1}_F \leq \int_E \mathds{1}_F \frac{f}{\alpha} \leq \frac{1}{\alpha}\int_E f \]
\end{proof}

\begin{cor}
If $t > 0$, $m(f^{-1}([\alpha, \infty])) \leq e^{-t\alpha}\int_E e^{tf}$.
\end{cor}
\begin{proof}
Holds since $m(f^{-1}([\alpha, \infty])) = m((e^{tf})^{-1}([\alpha, \infty]))$.
\end{proof}

\begin{cor}
If $\int_E f < \infty$ then $f(x)$ is finite a.e.
\end{cor}
\begin{proof}
$m(f^{-1}(\infty)) \leq m(f^{-1}([n, \infty])) \leq \frac{1}{n}\int_E f \to 0$.
\end{proof}

\subsection{Integral for general $f: E \to \cl{R}$}

\begin{defn}
We say that a measurable $f: E \to \cl{R}$ is integrable if $\int_E f_+ + \int_E f_- = \int_E |f| < \infty$. If $f$ is integrable then we define $\int_E f = \int_E f_+ - \int_E f_-$
\end{defn}

\begin{lemma}
Integrable functions form a vector space in that if $f, g$ are integrable functions then so is $f + g$ and if $f$ is integrable then so is $\lambda f$ for any $\lambda \in \R$.
\end{lemma}
\begin{proof}
Note that $\int_E |f + g| \leq \int_E |f| + \int_E |g| < \infty$ and $\int_E |\lambda f| = |\lambda|\int_E |f| < \infty$ which proves the lemma.
\end{proof}

\begin{exercise}
Check that $\int_E \alpha f + \beta g = \alpha\int_E f + \beta\int_E g$ for all integrable $f, g$ and constants $\alpha, \beta$.
\end{exercise}

\begin{exercise}
Check that Fatou's lemma holds if $f_n \geq 0$ a.e.
\end{exercise}

\begin{thm}[Dominated Convergence Theorem]
Let $f_n: E \to \cl{R}$ be a sequence of measurable functions that converge pointwise a.e. to some $f: E \to \cl{R}$. Assume that there is an integrable $g$ such that $|f_n| \leq g$ a.e. Then $\lim_{n \to \infty} \int_E f_n = \int_E f$.
\end{thm}
\begin{proof}
By assumptions, $g + f_n, g - f_n$ are nonnegative a.e. so we can apply Fatou's lemma. This gives us 
\[ \int_E g + \int_E f = \int_E g + f = \int_E \lim_{n \to \infty}(g + f_n) \leq \liminf_{n \to \infty} \int_E g + f_n = \int_E g + \liminf_{n \to \infty}\int_E f_n \]
Similarly, we have 
\[ \int_E g - \int_E f \leq \int_E g + \liminf_{n \to \infty} \int_E -f_n = \int_E g - \limsup_{n \to \infty} \int_E f_n \]
Since $\int_E g = \int_E |g| < \infty$, we get 
\[ \int_E f \leq \liminf_{n \to \infty} \int_E f_n \leq \limsup_{n \to \infty} \int_E f_n \leq \int_E f \]
Thus we have $\lim_{n \to \infty} \int_E f_n = \int_E f$.
\end{proof}

\end{document}